\section{Boolean reflection and bounded search}\label{sec:boolean-search}

\subsection{Boolean conventions}

Use numerical Booleans \(0,1\).  Define primitive-recursive operations
\[
 \neg b,\qquad b\wedge c,\qquad b\vee c,\qquad b\Rightarrow c,
\]
by fixed case distinctions, and normalize every Boolean-producing term so that an object-level range theorem proves \(b=0\vee b=1\).  Equality and order on natural numbers are primitive recursive; let
\[
 \mathsf{eq}(s,t),\quad \mathsf{lt}(s,t),\quad \mathsf{le}(s,t)
\]
denote their Boolean codes.  The basic reflection library proves, for example,
\[
 \mathsf{eq}(s,t)=1\ \leftrightarrow\ s=t,
 \qquad
 \mathsf{lt}(s,t)=1\ \leftrightarrow\ s<t.
\]
The right-hand relations are themselves expressed in the fixed arithmetic language, for instance using truncated subtraction or bounded witnesses.  Once proved, application modules use only the reflected form.

\subsection{Finite folds}

For a Boolean predicate term \(p(i,\vec x)\), define primitive-recursive folds
\begin{align*}
 \mathsf{All}(0,\vec x)&=1,
 &\mathsf{All}(S n,\vec x)&=\mathsf{All}(n,\vec x)\wedge p(n,\vec x),\\
 \mathsf{Any}(0,\vec x)&=0,
 &\mathsf{Any}(S n,\vec x)&=\mathsf{Any}(n,\vec x)\vee p(n,\vec x).
\end{align*}
These functions carry all bounded quantifier content in a numerical state.

\begin{theorem}[Finite-fold reflection]\label{thm:fold-reflection}
Assume \(p(i,\vec x)\) is Boolean-valued and reflects a fixed formula \(P(i,\vec x)\).  Then \iPRA\ derives
\begin{align*}
 \mathsf{All}(N,\vec x)=1
   &\leftrightarrow \forall i<N\;P(i,\vec x),\\
 \mathsf{Any}(N,\vec x)=1
   &\leftrightarrow \exists i<N\;P(i,\vec x).
\end{align*}
Moreover the induction used to establish the fold equations has a quantifier-free matrix.
\end{theorem}

\begin{proof}
The crucial point is not to induct on either displayed bounded-quantifier formula.  Compile a finite scan whose numerical state contains the two accumulator Booleans and, conditionally, the first true and first false indices.  Its update is primitive recursive.  A single quantifier-free invariant records Boolean range together with
\begin{align*}
 \mathsf{Any}(n)=1&\to t_n<n\wedge p(t_n)=1,\\
 \mathsf{All}(n)=0&\to f_n<n\wedge p(f_n)=0.
\end{align*}
The same induction proves the pointwise insertion and extraction equations
\begin{align*}
 \mathsf{All}(N)=1\wedge j<N&\to p(j)=1,\\
 p(j)=1\wedge j<N&\to \mathsf{Any}(N)=1.
\end{align*}
Every formula used here is quantifier-free: the candidate indices are stored terms, the bounds are reflected comparisons, and the finite prefix content is carried by numerical accumulator equations.  Equivalently, the stored-witness clauses may be discharged by two instances of the checked least-zero search backend, one applied to \(1\mathbin{\dot-}p\) and one to \(p\).

Now assemble the logical equivalences without further induction.  If \(\mathsf{All}(N)=1\), fix arbitrary \(j<N\); the first pointwise theorem gives \(p(j)=1\), and predicate reflection gives \(P(j)\).  Conversely assume \(\forall j<N\,P(j)\).  The Boolean range theorem gives \(\mathsf{All}(N)=0\) or \(1\).  In the zero case the stored false index \(f_N<N\) satisfies \(p(f_N)=0\), while the universal assumption and predicate reflection give \(p(f_N)=1\), a contradiction.  Hence \(\mathsf{All}(N)=1\).

If \(\mathsf{Any}(N)=1\), the stored true index \(t_N<N\) and reflection yield a bounded witness for \(P\).  Conversely, eliminate a supplied witness \(j<N\) with \(P(j)\), reflect it to \(p(j)=1\), and apply the insertion theorem.  Thus the only use of the \(\PRAAx\) induction rule is the quantifier-free scan invariant; the bounded quantifiers occur only in conclusions assembled afterward by the ordinary logical rules.
\end{proof}

\subsection{Structural Booleanization}

A \emph{fixed bounded formula} is generated from atomic equalities and inequalities by propositional connectives and quantifiers of the forms
\[
  \forall i<t(\vec x)\;A(i,\vec x),
  \qquad
  \exists i<t(\vec x)\;A(i,\vec x),
\]
where all terms are primitive recursive and the formula syntax is fixed at compilation time.

\begin{theorem}[Bounded Booleanization]\label{thm:booleanization}
For every fixed bounded formula \(A(\vec x)\), one can effectively construct a primitive-recursive Boolean term \(\chi_A(\vec x)\) and an object-level derivation
\[
  \vdash[\PRAAx] A(\vec x)\leftrightarrow \chi_A(\vec x)=1.
\]
The construction is syntax directed and uses no universal evaluator.
\end{theorem}

\begin{proof}
Proceed by structural recursion in the metatheory on the fixed formula \(A\).
\begin{itemize}
  \item For an atomic equation or comparison, use the corresponding reflected Boolean.
  \item For \(B\wedge C\), \(B\vee C\), and \(B\to C\), combine the recursively constructed Booleans with fixed Boolean operations and compose their reflection theorems in intuitionistic logic.  Implication is valid because the antecedent Boolean is decidable; its truth table is checked by the four Boolean cases.
  \item For a bounded universal or existential, recursively Booleanize the matrix and apply \cref{thm:fold-reflection} at the bound term.
\end{itemize}
The compiler constructs a dedicated term for this one syntax tree.  It does not code formulas and evaluate their codes inside the object theory.  Thus the construction is compatible with the impossibility of a total universal primitive-recursive evaluator.
\end{proof}

\subsection{Least bounded search}

Let \(p(i,\vec x)\) be a reflected Boolean predicate.  Define a search state \((f,w)\), where \(f\) is the found flag and \(w\) is the stored witness:
\[
 (f_0,w_0)=(0,0),
 \qquad
 (f_{i+1},w_{i+1})=
 \begin{cases}
   (1,i),&f_i=0\text{ and }p(i)=1,\\
   (f_i,w_i),&\text{otherwise.}
 \end{cases}
\]
The conditional is a primitive-recursive selector.

The useful invariant is not written with quantifiers.  Define Boolean prefix folds
\[
 E(i)=\mathsf{Any}_{j<i}p(j),
 \qquad
 L(i,w)=\mathsf{All}_{j<w}\neg p(j).
\]
Then store or recompute the following QF equalities:
\begin{align*}
 f_i&=E(i),\\
 f_i=1&\Rightarrow w_i<i,\\
 f_i=1&\Rightarrow p(w_i)=1,\\
 f_i=1&\Rightarrow L(i,w_i)=1,\\
 f_i=0&\Rightarrow w_i=0.
\end{align*}
Because \(E\) and \(L\) are terms, this is a quantifier-free formula.

\begin{theorem}[Least bounded search]\label{thm:least-search}
There are primitive-recursive terms \(\mathsf{found}(B,\vec x)\) and \(\mathsf{least}(B,\vec x)\) such that \iPRA\ derives:
\begin{enumerate}[label=(\alph*)]
  \item \(\mathsf{found}=1\leftrightarrow\exists i<B\;P(i,\vec x)\);
  \item if \(\mathsf{found}=1\), then \(\mathsf{least}<B\) and \(P(\mathsf{least},\vec x)\);
  \item if \(\mathsf{found}=1\) and \(j<\mathsf{least}\), then \(\neg P(j,\vec x)\);
  \item if no bounded witness exists, then \(\mathsf{found}=0\) and the witness output is the specified default.
\end{enumerate}
\end{theorem}

\begin{proof}
Encode the pair state by a fixed primitive-recursive pairing function and define its update by primitive recursion.  Prove the displayed QF invariant by \cref{lem:qfind}.  The base case follows from the empty-fold equations.  For the successor step, split on the two Boolean values of \(f_i\) and \(p(i)\).  If a witness was already found, the state is unchanged and every invariant clause persists.  If no witness was found and \(p(i)=0\), the empty status persists and the prefix fold remains false.  If no witness was found and \(p(i)=1\), the new witness is \(i\); validity is immediate, and minimality follows from the previous assertion \(E(i)=0\), reflected as the absence of an earlier witness.

Instantiate the invariant at \(B\).  Fold reflection turns \(f_B=E(B)\) into part (a).  The remaining clauses give validity and minimality.  The default statement is the final invariant clause.
\end{proof}


\subsection{Bounded-existential normal form and constructive Skolemization}

The preceding theorem does more than implement a search routine.  It gives a normal form for every bounded existential assertion.  Let \(P(i,\vec x)\) be a fixed bounded formula, let \(\chi_P\) be its reflected Boolean, and let \(B(\vec x)\) be an arbitrary primitive-recursive bound.  Write
\[
  e(\vec x)=\mathsf{found}(B(\vec x),\vec x),
  \qquad
  w(\vec x)=\mathsf{least}(B(\vec x),\vec x).
\]
Define the quantifier-free certificate matrix
\[
 \mathsf{LeastCert}_{P,B}(\vec x) \;\defeq\;
 e=1\;\wedge\;w<B\;\wedge\;\chi_P(w,\vec x)=1
 \;\wedge\;\mathsf{All}_{j<w}\neg\chi_P(j,\vec x)=1.
\]
Every symbol in this display is a term, so the matrix is quantifier-free even though it records bounded validity and leastness.

\begin{theorem}[Bounded-existential normal form]\label{thm:bounded-ex-normal}
For every fixed bounded formula \(P(i,\vec x)\) and primitive-recursive bound \(B(\vec x)\), one can construct terms \(e,w\) and object-level derivations of
\begin{align*}
 \exists i<B\;P(i,\vec x)
   &\leftrightarrow e(\vec x)=1,\\
 e(\vec x)=1
   &\leftrightarrow \mathsf{LeastCert}_{P,B}(\vec x).
\end{align*}
Consequently the bounded existential is interchangeable, by cut, with one quantifier-free certificate equation plus its fixed decoding lemmas.
\end{theorem}

\begin{proof}
Booleanize \(P\), form the bounded existential fold \(E(B)=\mathsf{Any}_{i<B}\chi_P(i)\), and apply \cref{thm:least-search}.  The search invariant proves \(e=E(B)\).  Fold reflection proves
\(E(B)=1\leftrightarrow\exists i<B\,P(i)\), giving the first equivalence.

For the forward implication of the second equivalence, the final search invariant supplies \(w<B\), \(\chi_P(w)=1\), and the prefix fold asserting that every earlier Boolean is zero.  Together with \(e=1\) these are exactly the certificate matrix.  The reverse implication is immediate by conjunction elimination.  No choice principle occurs: \(w\) is a primitive-recursive term defined for all inputs, and the validity theorem is conditional on its found flag.

To use the normal form in a larger proof, first replace an assumption
\(\exists i<B\,P(i)\) by \(e=1\) through the first equivalence.  Cut in the forward certificate theorem and reason with the explicit term \(w\).  Conversely, when a consumer produces the certificate, validity of \(w\) and existential introduction recover the bounded existential.  Since the loop invariant mentions only \(e,w\), Boolean folds, and equations, an induction that previously quantified over all earlier indices is now a legal quantifier-free induction.
\end{proof}

\begin{corollary}[Bounded witness extraction under a precondition]\label{cor:bounded-skolem}
Suppose \(H(\vec x)\) is any formula and iPRA derives
\[
 H(\vec x)\to\exists i<B(\vec x)\;P(i,\vec x).
\]
Then the same theory derives
\[
 H(\vec x)\to
 \bigl(w(\vec x)<B(\vec x)\wedge P(w(\vec x),\vec x)\bigr),
\]
where \(w\) is the fixed least-search term.  If a retained conclusion \(C(\vec x)\) follows uniformly from
\(i<B\wedge P(i,\vec x)\), then it follows from \(H\) after substituting \(w\) for \(i\).
\end{corollary}

\begin{proof}
Under \(H\), the first equivalence in \cref{thm:bounded-ex-normal} gives \(e=1\).  The search-validity clauses then give the bound and reflected predicate at \(w\); predicate reflection gives \(P(w)\).  Substitution of the term \(w\) into the consumer derivation and one cut yield the final statement.  Leastness is available when useful but is not required for ordinary witness extraction.
\end{proof}

This corollary is the constructive form of Skolemization that the framework uses.  It is valid only because the existential is bounded by an explicit primitive-recursive term and its matrix is decidable.  An unbounded existential or an assertion of mere finiteness does not supply such a term.

\subsection{Derived finite algorithms}

The same search theorem yields a library of reusable combinators.
\begin{itemize}
  \item \textbf{Bounded argmin.}  Search the least index attaining the least value in a finite table; comparison reflection supplies the predicate.
  \item \textbf{Duplicate-free enumeration.}  Scan a bounded box, append an item only if its verifier succeeds and no earlier appended item is equivalent.
  \item \textbf{Finite cover selection.}  Search a finite net for the first cell whose certified test covers the current query.
  \item \textbf{Pivot selection.}  Under a margin certificate asserting that some pivot exceeds a threshold, search the least such pivot.  Without the threshold, do not claim an exact pivot.
  \item \textbf{Finite path successor.}  Given a viability table, search the least compatible viable successor.
\end{itemize}
Each derived algorithm inherits a QF invariant because every finite universal or existential subcondition has first been replaced by a Boolean fold.
